\chapter{Codex CLI: A Practitioner's Reference}
\label{app:codex}

\begin{remark}[Learning Objectives]
After working through this appendix, the reader should be able to:
\begin{enumerate}
  \item Install Codex CLI, authenticate it, and select the appropriate approval
        mode for a given task.
  \item Explain the AGENTS.md file hierarchy and write project-level custom
        instructions for a quantitative finance codebase.
  \item Describe the three kernel-level sandboxing mechanisms (Seatbelt,
        Landlock, seccomp) and explain when to prefer Codex CLI's OS-enforced
        isolation over application-layer permission controls.
  \item Deploy a cloud-sandbox Codex task for processing sensitive financial
        documents without exposing data to the local development environment.
  \item Write and install a SKILL.md file that encodes a reusable finance
        research workflow.
  \item Compare the architectural trade-offs between Codex CLI and Claude Code
        and select the appropriate tool for a given research task.
  \item Map at least three quantitative-finance research tasks to concrete
        Codex CLI workflows, including an autonomous multi-hour data pipeline.
\end{enumerate}
\end{remark}

% ---------------------------------------------------------------
\section{What is Codex CLI}
\label{app:codex:overview}

\begin{context}
On 16 April 2025, OpenAI released Codex CLI as an open-source terminal agent.
The release was a deliberate competitive response: within days of Anthropic's
Claude Code reaching widespread adoption, OpenAI shipped a tool with a similar
surface---reads files, runs tests, edits code, commits changes---but with a
fundamentally different security architecture.  Where Claude Code enforces
safety at the application layer through programmable hooks and an allowlist,
Codex CLI enforces safety at the operating-system kernel layer using platform
sandboxing primitives.  By early 2026, approximately four million developers
were using Codex each week, and GPT-5.5---OpenAI's first fully agentic-first
retrained base model---had become the default engine behind every Codex surface
\citep{openai2026gpt55}.
\end{context}

\subsection{Overview and Positioning}
\label{app:codex:positioning}

Codex CLI is a terminal-based agentic coding assistant that reads a repository,
edits files, runs shell commands, and iterates on the results.  It is the
command-line surface of a broader \emph{Codex} platform that also includes a
cloud-hosted web interface, an IDE extension, a GitHub bot, and computer-use
capabilities.  All surfaces share the same underlying model and account context.

The most important architectural distinction from Claude Code
(Appendix~\ref{app:claude-code}) is the execution environment.  Claude Code
operates directly on the developer's local machine and file system; Codex CLI
can optionally delegate tasks to ephemeral cloud virtual machines pre-loaded
with the repository.  This cloud-sandbox mode provides reproducible, isolated
execution that is particularly attractive for sensitive financial data pipelines
where local environment contamination or credential leakage is a concern.

\subsection{Supported Models}
\label{app:codex:models}

As of mid-2026, Codex runs on the following model tiers:

\begin{itemize}
  \item \textbf{GPT-5.5} --- the default and recommended model; agentic-first
        training enables multi-hour autonomous sessions on real engineering tasks.
        Context window: 272K tokens (default); up to 1.05M with long-context mode
        enabled.
  \item \textbf{GPT-5.5 Pro} --- higher rate limits and priority routing for
        teams running parallel Codex agents on large pipelines.
  \item \textbf{GPT-4.1} --- lower cost option for high-frequency, shorter
        tasks such as linting, formatting, or schema validation.
\end{itemize}

Switch models with the \texttt{--model} flag or the \texttt{model} key in the
configuration file.  GPT-5.5's explicit agentic-first training means it is
significantly better than earlier models at sustaining multi-step task state
across tool calls---an important property for long financial data pipelines.

\subsection{Key Milestones}
\label{app:codex:milestones}

Table~\ref{tab:codex-milestones} summarises the major capability releases.

\begin{table}[ht]
\centering
\caption{Major Codex CLI capability releases (2025--2026).}
\label{tab:codex-milestones}
\begin{tabular}{lll}
\toprule
\textbf{Date} & \textbf{Feature} & \textbf{Notes} \\
\midrule
April 2025     & Codex CLI launch      & Open-source terminal agent; three approval modes \\
May 2025       & Cloud sandbox         & Ephemeral VM execution; pre-loaded repository \\
September 2025 & IDE extension         & VS Code and JetBrains integration \\
October 2025   & GitHub bot            & \texttt{@codex} in PR comments triggers a task \\
December 2025  & Skills support        & SKILL.md spec; shared with Claude Code \\
February 2026  & Computer use          & Screen reading; desktop app interaction \\
April 2026     & GPT-5.5 default       & Agentic-first model; 1.05M max context \\
\bottomrule
\end{tabular}
\end{table}

% ---------------------------------------------------------------
\section{Installation and Setup}
\label{app:codex:setup}

\subsection{Prerequisites and Authentication}
\label{app:codex:install}

Codex CLI requires Node.js 18 or later and is installed via \texttt{npm}:

\begin{minted}{bash}
npm install -g @openai/codex
\end{minted}

On first run, Codex opens a browser-based OAuth flow to authenticate against a
ChatGPT account.  For headless or CI environments, set the
\texttt{OPENAI\_API\_KEY} environment variable:

\begin{minted}{bash}
export OPENAI_API_KEY="sk-..."
codex                            # start interactive session
codex "analyse the factor model in src/models/ff5.py"
\end{minted}

To specify a model or run non-interactively:

\begin{minted}{bash}
codex --model gpt-5.5 --approval-mode full-access \
      "Refactor the CAPM implementation and add unit tests."
\end{minted}

\subsection{First-Run Configuration}
\label{app:codex:firstrun}

On first run, Codex creates a global configuration file at
\texttt{\textasciitilde/.codex/config.json}.  The most important top-level
keys are:

\begin{minted}{json}
{
  "model": "gpt-5.5",
  "approvalMode": "auto",
  "projectDocMaxBytes": 32768,
  "projectDocFallbackFilenames": ["AGENTS.md", "README.md"],
  "sandbox": {
    "mode": "local",
    "networkAccess": false
  }
}
\end{minted}

\texttt{projectDocMaxBytes} controls how much of the AGENTS.md file is read
into context on session start (default 32\,KB).  For large codebases with
detailed AGENTS.md conventions, increase this limit to ensure all conventions
are loaded.

% ---------------------------------------------------------------
\section{Core Interface and Approval Modes}
\label{app:codex:interface}

\subsection{CLI Commands and Navigation}
\label{app:codex:cli}

Codex CLI's built-in slash commands are the primary navigation layer.
Table~\ref{tab:codex-slash} lists the most commonly used commands.

\begin{table}[ht]
\centering
\caption{Built-in Codex CLI slash commands.}
\label{tab:codex-slash}
\begin{tabular}{ll}
\toprule
\textbf{Command} & \textbf{Effect} \\
\midrule
\texttt{/plan}        & Produce a plan without executing any changes \\
\texttt{/exec}        & Execute the current plan \\
\texttt{/review}      & Review all pending changes before committing \\
\texttt{/permissions} & Switch approval mode interactively \\
\texttt{/model}       & Switch the active model \\
\texttt{/context}     & Show current context usage and limits \\
\texttt{/sandbox}     & Toggle between local and cloud-sandbox mode \\
\texttt{/clear}       & Start a fresh session \\
\texttt{/help}        & List all available commands and skills \\
\bottomrule
\end{tabular}
\end{table}

The \texttt{/plan} → \texttt{/exec} → \texttt{/review} loop is the
recommended structured workflow for any task that modifies production data
or trading logic: plan in read-only mode, inspect the diff, then execute.

\subsection{The Three Approval Modes}
\label{app:codex:approval-modes}

Approval modes govern how much Codex can do without pausing for confirmation.
Unlike Claude Code's fine-grained allowlist/denylist syntax, Codex CLI offers
three coarse-grained modes:

\begin{description}
  \item[\texttt{read-only} (Suggest mode)] Codex can browse files and propose
        a plan but cannot write files or run shell commands until the user
        approves.  The safest mode for initial exploration of an unfamiliar
        codebase or dataset.

  \item[\texttt{auto} (Default)] Codex can read files, edit files, and run
        shell commands within the working directory.  It still pauses before
        accessing anything outside the working directory or performing network
        operations.  Equivalent to Claude Code's Edit mode with a constrained
        allowlist.

  \item[\texttt{full-access} (Danger mode)] Codex operates without
        confirmation across the entire machine, including network access.  Use
        sparingly and only in sandboxed environments.
\end{description}

Switch modes at any time with \texttt{/permissions} or by restarting with the
\texttt{--approval-mode} flag.

\begin{deepdive}
The coarse three-mode model is a deliberate design choice.  OpenAI's
reasoning: most developers find fine-grained allowlists difficult to maintain
correctly, so a well-enforced kernel-level boundary (see
Section~\ref{app:codex:sandboxing}) is a more reliable safety guarantee than
an application-layer list that can be misconfigured.  The trade-off is
flexibility: Claude Code's hook system can express policies like ``allow
\texttt{git commit} but deny \texttt{git push}'', whereas Codex cannot without
a custom AGENTS.md guard script.
\end{deepdive}

\subsection{Context Management}
\label{app:codex:context}

GPT-5.5's default context window is 272K tokens.  Long-context mode (up to
1.05M) is activated by setting \texttt{"longContext": true} in the
configuration file.  Key practical points:

\begin{itemize}
  \item Use \texttt{/context} to inspect current usage before starting a
        long pipeline; large repositories loaded in full can exhaust the
        default window before the task completes.
  \item Unlike Claude Code's \texttt{/compact} command, Codex does not offer
        manual compaction.  When the context limit is approached, Codex
        automatically summarises earlier turns.
  \item For pipelines that operate on many files sequentially (e.g., processing
        hundreds of earnings call transcripts), prefer the cloud-sandbox mode
        (Section~\ref{app:codex:cloud}) where each task starts with a fresh
        context.
\end{itemize}

% ---------------------------------------------------------------
\section{Configuration}
\label{app:codex:config}

\subsection{The AGENTS.md File Hierarchy}
\label{app:codex:agents-md}

\texttt{AGENTS.md} is Codex CLI's analogue of Claude Code's
\texttt{CLAUDE.md}.  It provides project-level custom instructions that Codex
reads at the start of every session.  The hierarchy follows the same layered
pattern:

\begin{enumerate}
  \item \textbf{Global} (\texttt{\textasciitilde/.codex/AGENTS.md}) ---
        personal preferences applied to all projects.
  \item \textbf{Project} (\texttt{AGENTS.md} at the repo root) --- project
        conventions; checked into version control.
  \item \textbf{Sub-directory} (any \texttt{AGENTS.md} in a nested folder) ---
        package-level or module-level overrides.
\end{enumerate}

Codex merges all three levels on startup.  The
\texttt{projectDocFallbackFilenames} configuration key allows additional files
(e.g., \texttt{README.md}) to serve as the project instruction source when
\texttt{AGENTS.md} is absent.

A representative AGENTS.md for a quantitative finance project:

\begin{minted}{markdown}
# Project: Equity Factor Research

## Conventions
- Always use log returns: `np.log(price / price.shift(1))`.
- Date columns must be `datetime64[ns]` with timezone `UTC`.
- Factor returns are stored in `data/factors/` as Parquet files.

## Testing
- Run `pytest tests/ -x -q` after every file edit.
- Coverage must remain above 80%.

## Data Access
- Raw price data: `data/raw/prices.parquet` (read-only).
- Bloomberg data requires the `blpapi` conda environment.
\end{minted}

\subsection{Configuration Reference}
\label{app:codex:config-ref}

Beyond the top-level keys shown in Section~\ref{app:codex:firstrun}, the
configuration file supports:

\begin{description}
  \item[\texttt{sandbox.mode}] \texttt{"local"} (default) or
        \texttt{"cloud"} (ephemeral VM).
  \item[\texttt{sandbox.networkAccess}] Boolean; default \texttt{false} in
        auto mode.  Set \texttt{true} only when the task requires fetching
        external data.
  \item[\texttt{longContext}] Boolean; enables GPT-5.5's 1.05M context window.
  \item[\texttt{historySize}] Number of past sessions stored in
        \texttt{\textasciitilde/.codex/history/}; default 50.
  \item[\texttt{shell}]  Shell used for command execution; default
        \texttt{"/bin/bash"} on macOS/Linux, \texttt{"powershell"} on Windows.
\end{description}

\subsection{Skills}
\label{app:codex:skills}

Codex CLI implements the same open agent skills specification as Claude Code.
A skill is a directory in \texttt{\textasciitilde/.agents/skills/} containing
a required \texttt{SKILL.md} plus optional scripts and reference assets:

\begin{minted}{text}
~/.agents/skills/
  fetch-factor-returns/
    SKILL.md           # name, description, invocation instructions
    scripts/
      download.py      # optional supporting script
\end{minted}

Codex loads a skill when the user invokes \texttt{/skill-name} or when the
task description semantically matches the skill's \texttt{description} field.
Because the skill format is identical to Claude Code's, a well-written
SKILL.md works in both tools without modification---a significant practical
advantage for teams that use both.

% ---------------------------------------------------------------
\section{Sandboxing and the Security Model}
\label{app:codex:sandboxing}

\begin{context}
The choice of where to enforce security boundaries has deep implications for
what can go wrong when an AI agent makes a mistake or is manipulated by
adversarial input.  Claude Code enforces safety at the application layer:
a misconfigured allowlist or a prompt-injection attack can potentially
cause the model to execute a command that the developer did not intend.
Codex CLI enforces safety at the operating-system kernel layer.  Even if the
model requests a forbidden operation, the kernel refuses it before any harm
occurs.  For financial institutions with strict data governance requirements,
this distinction matters.
\end{context}

\subsection{Local Sandboxing}
\label{app:codex:local-sandbox}

In local mode, Codex enforces sandboxing using platform-native kernel
primitives:

\begin{description}
  \item[macOS --- Apple Seatbelt] A mandatory access control framework built
        into the XNU kernel.  Codex's Seatbelt profile grants read access to
        the working directory and system libraries, write access to the working
        directory only, and no outbound network access (in auto mode).  The
        profile is applied via \texttt{sandbox-exec} and cannot be bypassed by
        the model or by malicious code running inside the session.

  \item[Linux --- Landlock] A Linux Security Module (available since kernel
        5.13) that implements a stackable, unprivileged access-control
        mechanism \citep{linux2022landlock}.  Codex uses Landlock to restrict
        file-system access to the working directory tree.

  \item[Linux --- seccomp] A kernel mechanism that filters system calls.
        Codex's seccomp profile blocks dangerous syscalls (e.g.,
        \texttt{execve} outside the working directory, raw socket creation)
        regardless of what the agent requests.
\end{description}

These three mechanisms work at layers below the process memory space, making
them resistant to prompt injection, dependency confusion attacks, and malicious
code execution via \texttt{eval}-style patterns.  Table~\ref{tab:sandbox-compare}
compares the two approaches.

\begin{table}[ht]
\centering
\caption{Security enforcement layer: Codex CLI vs Claude Code.}
\label{tab:sandbox-compare}
\begin{tabular}{lll}
\toprule
\textbf{Property} & \textbf{Codex CLI} & \textbf{Claude Code} \\
\midrule
Enforcement layer      & OS kernel (Seatbelt/Landlock/seccomp) & Application layer \\
Bypass risk            & Very low (kernel cannot be overridden) & Low (requires misconfigured allowlist) \\
Granularity            & Coarse (3 modes)  & Fine (per-tool patterns) \\
Custom policies        & Limited (AGENTS.md scripts) & Rich (allowlist + denylist syntax) \\
Network control        & Binary (on/off per mode) & Per-command pattern matching \\
Platform               & macOS, Linux      & All platforms \\
\bottomrule
\end{tabular}
\end{table}

\subsection{Cloud Sandbox Execution}
\label{app:codex:cloud}

Setting \texttt{"sandbox.mode": "cloud"} routes each Codex task to an
ephemeral virtual machine on OpenAI's infrastructure.  Codex pre-loads the VM
with the repository (fetched from GitHub or supplied as a tarball) and destroys
it after the task completes.  Benefits for financial workflows:

\begin{itemize}
  \item \textbf{Data isolation} --- sensitive datasets remain in the cloud VM
        and are not written to the developer's local disk.
  \item \textbf{Reproducibility} --- each task starts from a clean environment,
        eliminating dependency on local Python versions, conda environments, or
        system libraries.
  \item \textbf{Parallelism} --- multiple tasks can run simultaneously in
        separate VMs, enabling portfolio-wide analyses that would block a single
        local machine for hours.
  \item \textbf{Audit trail} --- VM execution logs are retained and exportable,
        supporting MiFID II and similar regulatory obligations around algorithmic
        decision-making.
\end{itemize}

Cloud sandbox mode is activated per-session:

\begin{minted}{bash}
codex --sandbox cloud \
      "Run the full factor backtest suite on the 2010--2024 dataset \
       and return a JSON summary of annualised Sharpe ratios."
\end{minted}

% ---------------------------------------------------------------
\section{Multi-Surface Architecture}
\label{app:codex:surfaces}

\subsection{CLI, IDE Extension, and ChatGPT}
\label{app:codex:surfaces-main}

The Codex name covers a family of surfaces that share a single model and
account context:

\begin{description}
  \item[Codex CLI] The terminal agent described in this appendix; the most
        capable surface for autonomous, multi-step workflows.

  \item[IDE Extension] Available for VS Code and JetBrains.  Provides inline
        diffs, a plan review pane, and \texttt{@}-mention for open files,
        analogous to Claude Code's IDE extensions.

  \item[ChatGPT Codex] The cloud-hosted version accessible through the
        ChatGPT interface.  Delegates tasks to cloud VMs; suitable for
        non-technical collaborators who need to trigger data analysis tasks
        without a terminal.
\end{description}

\subsection{GitHub Bot}
\label{app:codex:github-bot}

Mentioning \texttt{@codex} in a GitHub pull request comment or issue triggers
a Codex task that runs in a cloud sandbox against the repository at that
commit.  The bot can:

\begin{itemize}
  \item write unit tests for a new function described in the PR,
  \item propose fixes for failing CI checks, or
  \item analyse whether a change to a risk model breaks any existing
        assumptions documented in \texttt{AGENTS.md}.
\end{itemize}

Results are posted as a follow-up PR comment with a diff and a summary.

\subsection{Computer Use}
\label{app:codex:computer-use}

Codex's computer-use capability (released February 2026) allows the agent to
read the screen and interact with arbitrary desktop applications.  For finance
practitioners, this opens up workflows previously requiring manual data entry:

\begin{itemize}
  \item extracting structured data from Bloomberg terminal screens without an
        API licence,
  \item navigating legacy proprietary trading systems to retrieve historical
        fills or position data, and
  \item filling in regulatory reporting forms that do not expose a programmatic
        interface.
\end{itemize}

Computer use runs in a cloud VM to prevent screen-capture of sensitive local
displays.

% ---------------------------------------------------------------
\section{Cost and Token Management}
\label{app:codex:cost}

\subsection{Pricing Model}
\label{app:codex:pricing}

Codex CLI bills at standard OpenAI API rates for the selected model.
Table~\ref{tab:codex-token-costs} gives approximate consumption for
representative finance research tasks using GPT-5.5.

\begin{table}[ht]
\centering
\caption{Approximate token consumption for representative Codex CLI tasks (GPT-5.5).}
\label{tab:codex-token-costs}
\begin{tabular}{lrr}
\toprule
\textbf{Task} & \textbf{Input tokens} & \textbf{Output tokens} \\
\midrule
Explain a 500-line pricing library    & 8{,}000--15{,}000   & 1{,}000--3{,}000 \\
Classify one earnings call transcript & 10{,}000--20{,}000  & 1{,}000--2{,}000 \\
Refactor a factor model module        & 25{,}000--60{,}000  & 5{,}000--15{,}000 \\
Full backtest suite debug (multi-step) & 80{,}000--200{,}000 & 15{,}000--40{,}000 \\
\bottomrule
\end{tabular}
\end{table}

Cloud sandbox tasks additionally incur VM compute costs billed separately by
OpenAI; consult the pricing page for current rates.

\subsection{Cost-Reduction Strategies}
\label{app:codex:cost-reduction}

\begin{itemize}
  \item Use GPT-4.1 for high-frequency, low-complexity tasks (linting, schema
        validation, simple reformatting); reserve GPT-5.5 for tasks requiring
        sustained reasoning.
  \item Enable long-context mode only when the task genuinely requires it;
        longer contexts incur proportionally higher input costs.
  \item Use cloud sandbox parallelism for bulk batch tasks (e.g., classifying
        500 transcripts) rather than running them sequentially in a single
        long-context session.
  \item Store intermediate results to disk and restart sessions at natural
        checkpoints rather than allowing context to accumulate indefinitely.
\end{itemize}

% ---------------------------------------------------------------
\section{Codex CLI in Research and Finance Workflows}
\label{app:codex:finance}

\subsection{Use Cases for Quants and Researchers}
\label{app:codex:use-cases}

The following use cases are particularly well-suited to Codex CLI's
architecture:

\begin{description}
  \item[Large-document processing] OpenAI's own finance team used Codex to
        review 24{,}771 K-1 tax forms totalling 71{,}637 pages, accelerating
        the task by two weeks relative to the prior year \citep{openai2025codex}.
        The cloud sandbox mode kept all document content within OpenAI's
        infrastructure and off local developer machines.

  \item[Legacy code archaeology] Loading a proprietary options pricing library
        into a cloud sandbox and asking Codex to generate a call graph,
        annotate undocumented assumptions, and produce a test suite.  The
        kernel-level sandbox prevents the legacy code from accessing network
        resources or writing outside the working directory during execution.

  \item[Regulatory document parsing] Extracting structured data from SEC
        filings, ESMA reports, or FINMA circulars using the \texttt{/plan} →
        \texttt{/exec} structured loop, with \texttt{/review} before any
        database writes.

  \item[Automated backtesting] Delegating overnight backtest runs to cloud
        sandbox VMs: each strategy variant runs in its own VM with its own
        fresh environment, results are written to a shared S3 bucket, and a
        summary is posted to Slack on completion.

  \item[Computer-use data extraction] Using the computer-use capability to
        extract historical volatility surface data from a Bloomberg terminal
        screen directly into a Parquet file, bypassing the need for a Bloomberg
        API licence.
\end{description}

\subsection{Choosing Between Codex CLI and Claude Code}
\label{app:codex:comparison}

Both tools are capable general-purpose coding agents; the choice depends on
the task's security and ecosystem requirements.
Table~\ref{tab:codex-vs-claude} summarises the key decision dimensions.

\begin{table}[ht]
\centering
\caption{Decision guide: Codex CLI vs Claude Code (Appendix~\ref{app:claude-code}).}
\label{tab:codex-vs-claude}
\begin{tabular}{lll}
\toprule
\textbf{Dimension} & \textbf{Codex CLI} & \textbf{Claude Code} \\
\midrule
Security enforcement  & OS kernel (stronger guarantee)   & Application layer (more flexible) \\
Permission granularity & 3 coarse modes                  & Per-tool pattern matching \\
Cloud execution       & First-class (VM sandbox)         & Not built-in (CI/CD only) \\
Extensibility         & Skills + AGENTS.md               & Skills + hooks + MCP + plugins \\
Multi-agent           & Parallel cloud VMs               & Agent tool + Workflow + Agent Teams \\
GitHub integration    & \texttt{@codex} bot              & Via CI/CD non-interactive mode \\
Finance model access  & GPT-5.5 (OpenAI ecosystem)       & Claude Opus/Sonnet (Anthropic) \\
Context window        & 272K default / 1.05M max         & Up to 1M (Opus 4.8) \\
\bottomrule
\end{tabular}
\end{table}

\textbf{Choose Codex CLI when:}
\begin{itemize}
  \item kernel-level sandboxing is a regulatory or security requirement,
  \item the task involves processing sensitive documents that must not touch a
        local disk,
  \item your team already uses OpenAI's API and adding Anthropic would create
        vendor complexity, or
  \item you need cloud-parallel execution of many independent sub-tasks.
\end{itemize}

\textbf{Choose Claude Code when:}
\begin{itemize}
  \item fine-grained, per-command permission control is required (e.g.,
        allow \texttt{git commit} but deny \texttt{git push}),
  \item you need MCP server integration to connect directly to Bloomberg,
        FactSet, or an internal data API,
  \item the workflow requires programmable hooks that fire on specific tool
        events, or
  \item you are building a multi-agent system with a Workflow orchestration
        script and Agent Teams.
\end{itemize}

\subsection{Recommended Workflow Patterns}
\label{app:codex:patterns}

\begin{enumerate}
  \item \textbf{Plan before executing.}  Always start with \texttt{/plan} for
        any task that modifies a production dataset or trading strategy.
        Review the proposed diff carefully before typing \texttt{/exec}.

  \item \textbf{Use cloud sandbox for sensitive data.}  Configure
        \texttt{"sandbox.mode": "cloud"} in \texttt{\textasciitilde/.codex/config.json}
        for any project that handles NDA-covered data, client portfolios, or
        proprietary model weights.

  \item \textbf{Encode conventions in AGENTS.md.}  Every project should have
        a project-level \texttt{AGENTS.md} that specifies return conventions,
        data paths, testing requirements, and prohibited operations.  Codex
        reads this file at session start; consistent conventions eliminate an
        entire class of agent errors.

  \item \textbf{Separate concerns with skills.}  Build a SKILL.md for each
        recurring workflow step---factor construction, backtest execution, report
        generation---and invoke them by name.  Skills are portable: the same
        SKILL.md works in Claude Code, making it straightforward to run a
        workflow in either tool.

  \item \textbf{Commit checkpoints frequently.}  Have skills commit after each
        major step.  Small commits let you bisect regressions introduced by the
        agent without losing days of pipeline work.

  \item \textbf{Use \texttt{/review} before any destructive operation.}  For
        any task that writes to a database, deletes files, or posts to an
        external system, issue \texttt{/review} to inspect the full diff and
        confirm there are no unintended side effects.
\end{enumerate}

\begin{example}[End-to-end: regulatory filing extraction with Codex CLI]
\label{ex:codex-filing}
A compliance analyst needs to extract structured risk metrics from 200 ESMA
stress-test reports (PDF) and load them into a database.  Using Codex CLI in
cloud sandbox mode:

\begin{minted}{bash}
# Configure cloud sandbox for this project (data stays off local disk)
echo '{"sandbox": {"mode": "cloud"}}' > .codex/config.json

# Step 1: Plan the extraction pipeline (read-only, no changes yet)
codex --approval-mode read-only \
  "/plan Extract risk metrics (CET1 ratio, LCR, NSFR) from PDF files in
   data/esma_reports/. Output to data/structured/esma_metrics.parquet."

# Step 2: Review the plan, then execute
codex --approval-mode auto "/exec"

# Step 3: Validate schema and commit
codex "/review"
codex "Run tests/test_esma_schema.py and confirm all 200 reports parsed."
git add data/structured/ && git commit -m "feat: ESMA stress-test metrics"
\end{minted}

All PDF content is processed inside the cloud VM and never written to the
analyst's local disk.  The \texttt{/plan} → \texttt{/exec} → \texttt{/review}
loop provides an audit trail that satisfies internal model-governance
requirements.
\end{example}
